\documentclass[11pt]{article}
\usepackage{mathblatt}
% Präambel des Setzers (werkzeuge/setzer-praeambel.tex, wortgleich); dazu \gteil (Muster T2U)
\makeatletter
\setlength{\mbpfbe}{0mm}% keine Punkte (pruefung.md § 5: nur beim Prüfstein)
\newlength{\gEin}\setlength{\gEin}{0pt}
\let\gSchrift\relax
\renewcommand{\pfkreuzzeile}[1]{\par\noindent\hangindent3.2em\hangafter1\hspace*{1.6em}$\square$\ #1\par}
\newcommand{\gkreuz}[1]{$\square$\ #1\hspace{2em}}
\newcommand{\gkreuzzeile}[2]{\par\noindent\hangindent2.2em\hangafter1\mbox{#1\,$\square$}\ #2\par}% Buchstabe, Option: eine je Zeile (neben Skizze)
% links, was man ansieht, rechts, was man tut; Spaltengrenze überall an derselben Stelle (gemeinsam.md „Satz“)
% v0.17: Fußeintrag „Nr: Wert“ – kurz in einer Box (nie getrennt), lang umbrechbar mit der Nummer am Anfang
% (vorher \mbox: lange Lösungen liefen über den Rand, AS40 138 pt); robust, weil \fusshilfe in die .aux schreibt
\NewDocumentCommand{\gfuss}{m m}{\sbox0{#1:~#2}\ifdim\wd0<0.45\textwidth\mbox{#1:~#2}\else\mbox{#1:}\nobreak\ #2\fi}
\newsavebox{\gbox}\newlength{\gLinks}
% v0.10 (Lehrer 13.10.): keine starre Halbierung – die linke Spalte ist so breit wie die Figur (+ 6 mm), die
% Antwort steht direkt daneben; ist die Figur breiter als 0,62 Zeile, steht die Antwort darunter
\newcommand{\gzwei}[2]{\par\smallskip\noindent\sbox{\gbox}{#1}%
  \setlength{\gLinks}{\dimexpr\wd\gbox+6mm\relax}%
  \ifdim\gLinks>0.62\linewidth
    \usebox{\gbox}\par\smallskip\noindent\begin{minipage}{\linewidth}#2\end{minipage}%
  \else
    \begin{minipage}[t]{\gLinks}\vspace{0pt}\usebox{\gbox}\end{minipage}%
    \begin{minipage}[t]{\dimexpr\linewidth-\gLinks\relax}\vspace{0pt}\raggedright #2\end{minipage}%
  \fi\par}
% v0.9 (12.10., pruefung.md § 5 „Marke mit B, Bedingungen rechts“): drittes Argument von pfaufg = Bedingungen
% der Prüfung (v0.10: \pfohnetr = TR durchgestrichen, „FOR“), klein grau im rechten Rand auf der Grundlinie der Nummer, Zeilen mit \\
\newcommand{\gbedrand}[2]{\ifblank{#2}{}{\rlap{\hspace*{\dimexpr#1+2mm\relax}%
  \parbox[t]{15mm}{\raggedright\scriptsize\leavevmode\color{mbgrau}\strut#2}}}}
% Teil eines Bündels (a, b …): Marke links im Rand, Bedingungen rechts im Rand, auf der Zeile des Teils
\newcommand{\gteilrand}[2]{\llap{\mbpfrandmarke{#1}\hspace*{\mbpfnr}}\gbedrand{\linewidth}{#2}}
% Teil eines Bündels: Buchstabe in eigener Spalte, Inhalt eingerückt (\gEin = Einzug, damit die Spaltengrenze von
% \gzwei an derselben Stelle bleibt)
\newenvironment{gbteil}[3]{\par\smallskip\noindent\gteilrand{#2}{#3}\makebox[\mbpfteil][l]{\textbf{#1}}%
  \begin{minipage}[t]{\dimexpr\linewidth-\mbpfteil\relax}\setlength{\gEin}{\mbpfteil}\raggedright\strut\ignorespaces}%
 {\end{minipage}\par\addvspace{6pt}}% v0.18 (PG80 Nr. 1): Abstand, damit die Linie von a) nicht an b) rutscht
% Grafik höchstens so breit wie die Spalte (Bündel nebeneinander, Koordinatensystem ohne Karo)
\newsavebox{\gpbox}
\newcommand{\gpassend}[1]{\sbox{\gpbox}{#1}\ifdim\wd\gpbox>\linewidth\typeout{GPASSEND: verkleinert \the\wd\gpbox\space auf \the\linewidth}\resizebox{\linewidth}{!}{\usebox{\gpbox}}\else\usebox{\gpbox}\fi}
% Skizzenfeld (v0.10, Lehrer 13.10.): Karo ab dem linken Rand über die Zeile, „Skizze“ klein grau oben links im Feld
% (wie \pfskizzenfeld der Vorlage); der rechte Rand bleibt den Bedingungen
\newcommand{\gskizze}[1]{\par\smallskip\noindent\begin{tikzpicture}
  \clip (0,0) rectangle ({\linewidth-1mm},{#1*10mm});
  \draw[black!16,line width=0.3pt,step=5mm] (0,0) grid ({\linewidth},{#1*10mm});
  \node[anchor=north west,inner sep=2pt,font=\scriptsize,text=mbgrau,fill=white] at (0,{#1*10mm}) {Skizze};
  \end{tikzpicture}\par}
\RenewDocumentEnvironment{pfaufg}{m m m}{%
  \begin{lrbox}{\mbaufgabebox}\begin{minipage}[t]{\linewidth}%
  \gSchrift\noindent\hspace*{\gEin}\mbpfrandmarke{#1}\gbedrand{\linewidth-\gEin}{#3}%
  \begin{minipage}[t]{\mbpfnr}\strut\textbf{#2}\end{minipage}%
  \begin{minipage}[t]{\dimexpr\linewidth-\gEin-\mbpfnr\relax}\raggedright\strut\ignorespaces}%
 {\end{minipage}%
  \end{minipage}\end{lrbox}%
  \setlength{\mbaufgabehoehe}{\dimexpr\ht\mbaufgabebox+\dp\mbaufgabebox\relax}%
  \ifdefempty{\gkopfnach}{\par\addvspace{5pt}\Needspace*{\mbaufgabehoehe}}{%
    % Stufenkopf, der mit dieser (ersten) Aufgabe zusammenbleibt: Platz für beide verlangen, dann den Kopf setzen
    % (12.10., Quadratische: ein Kopf stand allein unten, die Aufgabe mit Koordinatensystem auf der nächsten Seite)
    % v0.17 (AS40 „Terme deuten“ allein unten): hinter dem Kopf steht das Sprungziel (\special, ein Knoten, nach
    % dem die Seite brechen durfte) – darum \nobreak nach dem Kopf; Vorhalt 56 pt = Kopf mit Linie und Abständen
    \par\Needspace*{\dimexpr\mbaufgabehoehe+56pt\relax}\gkopfnach\par\nobreak\gdef\gkopfnach{}\addvspace{5pt}}%
  \noindent\usebox{\mbaufgabebox}\par\addvspace{5pt}}
% v0.18 (plan.md § 10, 15.10. l): Abschnittsbuchstabe A, B, C vor dem Namen im Kopf – als Vorsatz, damit der
% Name im TeX gleich bleibt (Bestellkarte, Sprungziele lesen ihn)
\def\gbstakt{}\newcommand{\gbstvor}[1]{\gdef\gbstakt{#1}}
\NewCommandCopy\gpfstufekopfalt\pfstufekopf
\RenewDocumentCommand{\pfstufekopf}{o m m}{\IfValueTF{#1}{\gpfstufekopfalt[#1]}{\gpfstufekopfalt}%
  {\ifx\gbstakt\empty\else\gbstakt\enspace\fi #2}{#3}\gdef\gbstakt{}}
% v0.18 (16.10. e): Begriff grau, ohne Nummer, ohne Rahmen über der ersten Aufgabe der Stufe
\newcommand{\gbegriff}[1]{\par\nobreak\noindent\hspace*{\mbpfnr}\parbox[t]{\dimexpr\linewidth-\mbpfnr\relax}{\raggedright\small\color{mbgrau}#1}\par\nobreak\addvspace{3pt}}
% v0.18 (16.10. g): Lösungsblatt – Kopf mit kräftigem Strich, Abschnitt klein grau als Zwischenzeile
\newcommand{\glkopf}[1]{\noindent{\large\bfseries Lösungen \textperiodcentered{} #1}\par\vspace{3pt}\hrule height 1.2pt\par\medskip}
\newcommand{\glzwischen}[1]{\par\addvspace{8pt}\Needspace*{4\baselineskip}\noindent{\footnotesize\color{mbgrau}#1}\par\nobreak\vspace{1pt}}
% v0.18 (15.10. o, w): Blattende – Kern ohne Rahmen, dann graue Zeile „Mehr: n Aufgaben“
\newcommand{\gkern}[1]{\par\addvspace{14pt}\Needspace*{6\baselineskip}\noindent{\bfseries Kern}\par\nobreak\smallskip
  {\small\raggedright #1\par}}
\newcommand{\gmehr}[1]{\par\addvspace{10pt}\noindent{\footnotesize\color{mbgrau}Mehr: #1 Aufgaben}\par}
% Stufenkopf vormerken: \gkopfvor{<Kopf>} – die nächste pfaufg setzt ihn, sobald beide zusammen Platz haben
\def\gkopfnach{}
\newcommand{\gkopfvor}[1]{\gdef\gkopfnach{#1}}
% Teilaufgabe mit Skizze links in gzwei: Buchstabe in eigener Spalte, darunter Text und Skizze
\newcommand{\gteil}[2]{\makebox[\mbpfteil][l]{\textbf{#1}}\parbox[t]{\dimexpr0.5\textwidth-\gEin-\mbpfnr-5mm-\mbpfteil\relax}{\raggedright #2}}
% Stufenkopf der Vorlage (mathblatt.sty 2026-10-12: Linie, Abstand, 13 pt, \pfjahre) ohne ihren Vorhalt von
% 14 Zeilen: der Setzer fordert vor dem Kopf 0,16 Seitenhöhe (Kopf + erste Aufgabe), mehr nicht
\setlength{\mbpfstufevorhalt}{0pt}
\makeatother
\newcommand{\karo}[1]{\par\smallskip\noindent\begin{tikzpicture}
  \clip (0,0) rectangle ({\linewidth-1mm},{#1*10mm});
  \draw[black!16,line width=0.3pt,step=5mm] (0,0) grid ({\linewidth},{#1*10mm});
  \end{tikzpicture}\par}
\newcommand{\kk}[1]{$\square$\,#1\hspace{0.9em}}
\newcommand{\linie}[1]{\,{\color{black!45}\rule[-1pt]{#1}{0.4pt}}\,}
\newcommand{\pkt}{\textperiodcentered{}}
% v0.14 (14.10.): Stufenkopf mit fünf Punkten (Option = Jahre für die Bestellkarte) oder „zuletzt 19“
\NewDocumentCommand{\gpunkte}{o m}{#2}
\newcommand{\gzuletzt}[1]{zuletzt #1}
% v0.7 (12.10.): Schülerfrage unter dem Stufenkopf; Ankreuzoption als Bild (Buchstabe und Kästchen darunter)
\newcommand{\gfrage}[1]{\par\nobreak\noindent\hspace*{\mbpfnr}{\small\itshape #1}\par\nobreak\addvspace{2pt}}
\newlength{\gbildbreite}\setlength{\gbildbreite}{0.32\linewidth}% Setzer setzt sie je Reihe (vier Bilder: 0,24)
\newcommand{\gbildoption}[3]{\begin{minipage}[t][#3][b]{\gbildbreite}\centering #2\par\smallskip\mbox{#1\,$\square$}\end{minipage}}
% Lineare (12.10.): Bild mit Buchstabe ohne Kästchen (Diagramm zuordnen); Aussage mit „wahr / falsch“ zum Ankreuzen
\newcommand{\gbildname}[3]{\begin{minipage}[t][#3][b]{\gbildbreite}\centering #2\par\smallskip\mbox{\textbf{#1}}\end{minipage}}
\newcommand{\gwahrfalsch}[2]{\par\noindent\hangindent2.2em\hangafter1\makebox[2.2em][l]{(#1)}#2\hfill\mbox{wahr $\square$\quad falsch $\square$}\par}
\begin{document}

\pfheftstil
\pflaufkopf[P10 \pkt{} Prüfung]{Lineare Funktionen \pkt{} Lösungen}{PA10-7}\pfseitevon
\glkopf{Lineare Funktionen}
\small

\glzwischen{A~Gerade zeichnen}
\noindent{\footnotesize\color{mbgrau}Wo steht die Gleichung der linearen Funktion mit m und n? Formelsammlung, S. 31}\par
\begin{pfloesungb}{}
\lzb{1.}{}{a) und c)}{\begin{tabular}[t]{@{}>{\raggedleft\arraybackslash}p{0.34\linewidth}@{\,}>{\raggedright\arraybackslash}p{0.6\linewidth}@{}}\multicolumn{2}{@{}>{\raggedright\arraybackslash}p{0.97\linewidth}@{}}{linear sind a) und c): a) m = $-$2, n = 3 (die Zahl vor x ist m, auch hinten)}\\\multicolumn{2}{@{}>{\raggedright\arraybackslash}p{0.97\linewidth}@{}}{c) m = 0,5, n = 0}\\\multicolumn{2}{@{}>{\raggedright\arraybackslash}p{0.97\linewidth}@{}}{x² gibt keine Gerade}\\\end{tabular}}
\end{pfloesungb}
\begin{pfloesungb}{}
\lzb{2.}{}{die einzige Gerade; B Hyperbelast, C Parabel, D kein Funktionsgraph}{\begin{tabular}[t]{@{}>{\raggedleft\arraybackslash}p{0.34\linewidth}@{\,}>{\raggedright\arraybackslash}p{0.6\linewidth}@{}}\multicolumn{2}{@{}>{\raggedright\arraybackslash}p{0.97\linewidth}@{}}{die einzige Gerade}\\\multicolumn{2}{@{}>{\raggedright\arraybackslash}p{0.97\linewidth}@{}}{B Hyperbelast, C Parabel, D kein Funktionsgraph}\\\end{tabular}}
\end{pfloesungb}
\begin{pfloesungb}{}
\lzb{3.}{}{Punkt (0|1)}{\begin{tabular}[t]{@{}>{\raggedleft\arraybackslash}p{0.34\linewidth}@{\,}>{\raggedright\arraybackslash}p{0.6\linewidth}@{}}n & = 1 aus y = 1,5x + 1\\\multicolumn{2}{@{}>{\raggedright\arraybackslash}p{0.97\linewidth}@{}}{Punkt (0|1)}\\\end{tabular}}
\end{pfloesungb}
\begin{pfloesungb}{}
\lzb{4.}{}{Gerade durch (0|3) und (1|5)}{\begin{tabular}[t]{@{}>{\raggedleft\arraybackslash}p{0.34\linewidth}@{\,}>{\raggedright\arraybackslash}p{0.6\linewidth}@{}}\multicolumn{2}{@{}>{\raggedright\arraybackslash}p{0.97\linewidth}@{}}{Startpunkt (0|3), Steigungsdreieck 1 nach rechts, 2 nach oben}\\\multicolumn{2}{@{}>{\raggedright\arraybackslash}p{0.97\linewidth}@{}}{Gerade durch (0|3) und (1|5)}\\\end{tabular}}
\end{pfloesungb}

\glzwischen{B~Gleichung ablesen}
\begin{pfloesungb}{}
\lzb{5.}{a)}{h}{\begin{tabular}[t]{@{}>{\raggedleft\arraybackslash}p{0.34\linewidth}@{\,}>{\raggedright\arraybackslash}p{0.6\linewidth}@{}}\multicolumn{2}{@{}>{\raggedright\arraybackslash}p{0.97\linewidth}@{}}{Steigung $-$3: f}\\\multicolumn{2}{@{}>{\raggedright\arraybackslash}p{0.97\linewidth}@{}}{parallel zur x-Achse: h}\\\end{tabular}}
\end{pfloesungb}
\begin{pfloesungb}{}
\lzb{}{b)}{y = $-$1}{\begin{tabular}[t]{@{}>{\raggedleft\arraybackslash}p{0.34\linewidth}@{\,}>{\raggedright\arraybackslash}p{0.6\linewidth}@{}}f: y & = $-$3x + 3\\g: y & = x + 3\\h: y & = $-$1\\\end{tabular}}
\end{pfloesungb}

\glzwischen{C~Gleichung aus zwei Punkten}
\begin{pfloesungb}{}
\lzb{6.}{}{Gerade durch A und B; Aussage 1 falsch, Aussage 2 wahr; y = $-$1,5x + 2}{\begin{tabular}[t]{@{}>{\raggedleft\arraybackslash}p{0.34\linewidth}@{\,}>{\raggedright\arraybackslash}p{0.6\linewidth}@{}}\multicolumn{2}{@{}>{\raggedright\arraybackslash}p{0.97\linewidth}@{}}{A und B eintragen}\\\multicolumn{2}{@{}>{\raggedright\arraybackslash}p{0.97\linewidth}@{}}{Gerade durch A und B}\\m < 0, n & = 2\\\multicolumn{2}{@{}>{\raggedright\arraybackslash}p{0.97\linewidth}@{}}{Aussage 1 falsch, Aussage 2 wahr}\\m & = $\tfrac{- 2{,}5 - 5}{3 - ( - 2)}$ = $-$1,5\\n & = 5 $-$ ($-$1,5) · ($-$2) = 2\\y & = $-$1,5x + 2\\\end{tabular}}
\end{pfloesungb}

\glzwischen{D~Nullstelle gesucht}
\begin{pfloesungb}{}
\lzb{7.}{}{Gerade durch (0|5) und (1|3); x = 2,5}{\begin{tabular}[t]{@{}>{\raggedleft\arraybackslash}p{0.34\linewidth}@{\,}>{\raggedright\arraybackslash}p{0.6\linewidth}@{}}\multicolumn{2}{@{}>{\raggedright\arraybackslash}p{0.97\linewidth}@{}}{Startpunkt (0|5), 1 nach rechts, 2 nach unten}\\\multicolumn{2}{@{}>{\raggedright\arraybackslash}p{0.97\linewidth}@{}}{Gerade durch (0|5) und (1|3)}\\y & = 0: $-$2x + 5 = 0\\x & = 2,5\\\end{tabular}}
\end{pfloesungb}

\glzwischen{E~Punktprobe}
\begin{pfloesungb}{}
\lzb{8.}{}{$(3\,|\,-12)$}{\begin{tabular}[t]{@{}>{\raggedleft\arraybackslash}p{0.34\linewidth}@{\,}>{\raggedright\arraybackslash}p{0.6\linewidth}@{}}$-$12 $\neq$ f(3) & = $-$5 · 3 + 2 = $-$13\\\end{tabular}}
\end{pfloesungb}
\begin{pfloesungb}{}
\lzb{9.}{}{P liegt auf dem Graphen}{\begin{tabular}[t]{@{}>{\raggedleft\arraybackslash}p{0.34\linewidth}@{\,}>{\raggedright\arraybackslash}p{0.6\linewidth}@{}}f($-$2) & = $-$3 · ($-$2) + 1 = 7 = y-Wert von P\\\multicolumn{2}{@{}>{\raggedright\arraybackslash}p{0.97\linewidth}@{}}{P liegt auf dem Graphen}\\\end{tabular}}
\end{pfloesungb}

\glzwischen{F~Gleichung aus Text}
\begin{pfloesungb}{}
\lzb{10.}{}{a) 8 ist n, 0,09 ist m; b) 25 ist n, $-$2 ist m}{\begin{tabular}[t]{@{}>{\raggedleft\arraybackslash}p{0.34\linewidth}@{\,}>{\raggedright\arraybackslash}p{0.6\linewidth}@{}}\multicolumn{2}{@{}>{\raggedright\arraybackslash}p{0.97\linewidth}@{}}{n steht am Anfang (x = 0), m ist das, was je Monat, Minute oder Stunde dazukommt oder weggeht}\\\end{tabular}}
\end{pfloesungb}
\begin{pfloesungb}{}
\lzb{11.}{a)}{Maxi = 1648 €}{\begin{tabular}[t]{@{}>{\raggedleft\arraybackslash}p{0.34\linewidth}@{\,}>{\raggedright\arraybackslash}p{0.6\linewidth}@{}}Maxi & = 1 360 + 12 · 24 = 1648 €\\\end{tabular}}
\end{pfloesungb}
\begin{pfloesungb}{}
\lzb{}{b)}{y = 65x + 820; 16 Monate}{\begin{tabular}[t]{@{}>{\raggedleft\arraybackslash}p{0.34\linewidth}@{\,}>{\raggedright\arraybackslash}p{0.6\linewidth}@{}}y & = 65x + 820\\\multicolumn{2}{@{}>{\raggedright\arraybackslash}p{0.97\linewidth}@{}}{65x + 820 $\ge$ 1800}\\\multicolumn{2}{@{}>{\raggedright\arraybackslash}p{0.97\linewidth}@{}}{x $\ge$ $\tfrac{980}{65}$ $\approx$ 15,1}\\\multicolumn{2}{@{}>{\raggedright\arraybackslash}p{0.97\linewidth}@{}}{16 Monate}\\\end{tabular}}
\end{pfloesungb}
\begin{pfloesungb}{}
\lzb{12.}{a)}{Angebot 1 = C; Angebot 2 = B}{\begin{tabular}[t]{@{}>{\raggedleft\arraybackslash}p{0.34\linewidth}@{\,}>{\raggedright\arraybackslash}p{0.6\linewidth}@{}}\multicolumn{2}{@{}>{\raggedright\arraybackslash}p{0.97\linewidth}@{}}{erst fest (320 km frei), dann gleichmäßig steigend}\\Angebot 1 & = C\\\multicolumn{2}{@{}>{\raggedright\arraybackslash}p{0.97\linewidth}@{}}{Grundbetrag 115 €, dann gleichmäßig steigend}\\Angebot 2 & = B\\\end{tabular}}
\end{pfloesungb}
\begin{pfloesungb}{}
\lzb{}{b)}{Angebot 1 = 170,50 €; Angebot 2 = 151,00 €; Angebot 2 günstiger; y = 0,08x + 115}{\begin{tabular}[t]{@{}>{\raggedleft\arraybackslash}p{0.34\linewidth}@{\,}>{\raggedright\arraybackslash}p{0.6\linewidth}@{}}Angebot 1 & = 5 · 25 + (450 $-$ 320) · 0,35 = 170,50 €\\Angebot 2 & = 115 + 450 · 0,08 = 151,00 €\\\multicolumn{2}{@{}>{\raggedright\arraybackslash}p{0.97\linewidth}@{}}{151,00 € < 170,50 €}\\\multicolumn{2}{@{}>{\raggedright\arraybackslash}p{0.97\linewidth}@{}}{Angebot 2 günstiger}\\y & = 0,08x + 115\\\end{tabular}}
\end{pfloesungb}

\glzwischen{G~Gerade und Parabel}
\begin{pfloesungb}{}
\lzb{13.}{}{$x = -3$}{\begin{tabular}[t]{@{}>{\raggedleft\arraybackslash}p{0.34\linewidth}@{\,}>{\raggedright\arraybackslash}p{0.6\linewidth}@{}}$-$3 · ($-$3 + 4) & = $-$3 · 1 = $-$3\\\end{tabular}}
\end{pfloesungb}
\begin{pfloesungb}{}
\lzb{14.}{}{Gerade durch (0|1) und (1|$-$2)}{\begin{tabular}[t]{@{}>{\raggedleft\arraybackslash}p{0.34\linewidth}@{\,}>{\raggedright\arraybackslash}p{0.6\linewidth}@{}}\multicolumn{2}{@{}>{\raggedright\arraybackslash}p{0.97\linewidth}@{}}{Startpunkt (0|1), Steigungsdreieck 1 nach rechts, 3 nach unten}\\\multicolumn{2}{@{}>{\raggedright\arraybackslash}p{0.97\linewidth}@{}}{Gerade durch (0|1) und (1|$-$2)}\\\end{tabular}}
\end{pfloesungb}
\begin{pfloesungb}{}
\lzb{15.}{a)}{Nein}{\begin{tabular}[t]{@{}>{\raggedleft\arraybackslash}p{0.34\linewidth}@{\,}>{\raggedright\arraybackslash}p{0.6\linewidth}@{}}Nein, 5 $\neq$ p($-$3) & = ($-$3)² $-$ 5 = 4\\\end{tabular}}
\end{pfloesungb}
\begin{pfloesungb}{}
\lzb{}{b)}{S$_1$(4|11); S$_2$($-$2|$-$1)}{\begin{tabular}[t]{@{}>{\raggedleft\arraybackslash}p{0.34\linewidth}@{\,}>{\raggedright\arraybackslash}p{0.6\linewidth}@{}}x² $-$ 5 & = 2x + 3\\x² $-$ 2x $-$ 8 & = 0\\x$_1$ & = 4, x$_2$ = $-$2\\\multicolumn{2}{@{}>{\raggedright\arraybackslash}p{0.97\linewidth}@{}}{S$_1$(4|11)}\\\multicolumn{2}{@{}>{\raggedright\arraybackslash}p{0.97\linewidth}@{}}{S$_2$($-$2|$-$1)}\\\end{tabular}}
\end{pfloesungb}
\begin{pfloesungb}{}
\lzb{16.}{}{Gerade durch (0|$-$2) und (1|1); S(2|4); Aussage 2 wahr, Aussagen 1 und 3 falsch}{\begin{tabular}[t]{@{}>{\raggedleft\arraybackslash}p{0.34\linewidth}@{\,}>{\raggedright\arraybackslash}p{0.6\linewidth}@{}}\multicolumn{2}{@{}>{\raggedright\arraybackslash}p{0.97\linewidth}@{}}{Startpunkt (0|$-$2), 1 nach rechts, 3 nach oben}\\\multicolumn{2}{@{}>{\raggedright\arraybackslash}p{0.97\linewidth}@{}}{Gerade durch (0|$-$2) und (1|1)}\\\multicolumn{2}{@{}>{\raggedright\arraybackslash}p{0.97\linewidth}@{}}{Schnitt im Bild ablesen}\\\multicolumn{2}{@{}>{\raggedright\arraybackslash}p{0.97\linewidth}@{}}{S(2|4)}\\m & = 3 > 0, zwei Schnittpunkte im Bild, y-Achse bei (0|$-$2)\\\multicolumn{2}{@{}>{\raggedright\arraybackslash}p{0.97\linewidth}@{}}{Aussage 2 wahr, Aussagen 1 und 3 falsch}\\\end{tabular}}
\end{pfloesungb}
\begin{pfloesungb}{}
\lzb{17.}{}{S$_1$(1|3); S$_2$(5|$-$5)}{\begin{tabular}[t]{@{}>{\raggedleft\arraybackslash}p{0.34\linewidth}@{\,}>{\raggedright\arraybackslash}p{0.6\linewidth}@{}}(x $-$ 4)² $-$ 6 & = $-$2x + 5\\x² $-$ 6x + 5 & = 0\\x$_1$ & = 1, x$_2$ = 5\\\multicolumn{2}{@{}>{\raggedright\arraybackslash}p{0.97\linewidth}@{}}{S$_1$(1|3)}\\\multicolumn{2}{@{}>{\raggedright\arraybackslash}p{0.97\linewidth}@{}}{S$_2$(5|$-$5)}\\\end{tabular}}
\end{pfloesungb}
\end{document}